% TOP_PROBLEM: {"id": 30006955, "title": "Lusztig's P1-P15 Conjectures for Bounded Weighted Coxeter Groups", "category_id": 4, "category": {"id": 4, "name": "algebra", "display_name": "Algebra", "description": "Group theory, ring theory, field theory, and algebraic structures.", "slug": "algebra", "order_index": 4, "created_at": "2026-07-31T15:26:25.671Z"}, "status": "open"}
% FIRST_POSTED: 2026-09-27
% !TeX program = lualatex
% ATTEMPT_STATUS: unresolved
% Model: GPT-6 Astra Ultra; continuation dated 27 September 2026.
\documentclass[11pt]{article}
\usepackage[margin=25mm]{geometry}
\usepackage{fontspec}
\setmainfont{Latin Modern Roman}
\usepackage{amsmath,amssymb,mathtools}
\usepackage{unicode-math}
\setmathfont{Latin Modern Math}
\usepackage{xurl,hyperref}
\hypersetup{colorlinks=true,urlcolor=blue}
\setcounter{secnumdepth}{0}
\setlength{\parindent}{0pt}
\setlength{\parskip}{5pt}
\setlength{\emergencystretch}{3em}
\begin{document}
\section*{\texorpdfstring{Lusztig's P1-P15 Conjectures for Bounded Weighted Coxeter Groups}{Lusztig's P1-P15 Conjectures for Bounded Weighted Coxeter Groups}}\label{30006955}
\subsection{Definitions and mathematical statement}
Let $(W,S)$ be a Coxeter system with $S$ finite and length $\ell$. A positive weight function $L:W\to{\mathbb{Z}}$ satisfies $L(xy)=L(x)+L(y)$ whenever lengths add, and $L(s)>0$ for $s\in S$. Its unequal-parameter Hecke algebra over ${\mathbb{Z}}[v,v^{-1}]$ has standard basis $T_w$, the Coxeter braid relations, and quadratic relations
\[
 (T_s-v^{L(s)})(T_s+v^{-L(s)})=0.
\]
Writing $T_xT_y=\sum_zf_{x,y,z}T_z$, assume there is a uniform finite upper bound on the $v$-degrees of every $f_{x,y,z}$. This is precisely the boundedness hypothesis in Section 13.2 of the fixed 2014 revision.

Prove the conjunction of all fifteen statements P1--P15 in Lusztig's Conjectures 14.2 of that revision. These assertions concern the Kazhdan--Lusztig basis, its structure constants, the $a$-function, distinguished elements, and cells. Their full indexed identities and auxiliary definitions are incorporated by that exact reference rather than replaced by a weaker cell-structure slogan. Boundedness is assumed, not an additional conclusion to establish.
\par\smallskip\textit{Formulation and concept references:} \href{https://arxiv.org/pdf/math/0208154v2}{[S1]}.

\subsection{Short English statement}
Do all fifteen of Lusztig's specified unequal-parameter Hecke-algebra assertions hold whenever the weighted Coxeter group satisfies the stated boundedness hypothesis?
\subsection{Research attempt}

\paragraph{Fixed conventions and source audit (27 September 2026).}
The target is the conjunction in the specified 2014 revision [S1],
including its two-variable identity P15. We use the normalization of
the quadratic relation printed above, and write $C_w$ for the
bar-invariant Kazhdan--Lusztig element congruent to $T_w$ modulo
strictly negative powers of $v$ in lower Bruhat terms. The symbol
$C_w$ here denotes the source's $c_w$, not a differently normalized
primed basis. This convention affects both degrees and signs.

The inspected recent literature still distinguishes general unequal
parameters from solved classes. Xie's primary paper [E1] proves all
fifteen properties for Coxeter groups with complete graph and for
right-angled Coxeter groups. The elementary commuting case treated
below is contained in that established theory. Its value here is an
explicit proof and a fully reproducible certificate, not new generality.
A 2026 primary paper [E2] likewise distinguishes general unequal
parameters from its affine equal-parameter setting. Boundedness is
already a hypothesis of this entry; establishing a degree bound alone
would not answer the fifteen-property question.

The approach is to compute everything for a finite product of rank-one
systems, check each indexed property including P15, and then identify
the first obstruction when braid interactions are introduced.

\paragraph{Definitions used in the computation.}
Let $A=\mathbb Z[v,v^{-1}]$ and write
\[
 C_xC_y=\sum_z h_{x,y,z}C_z,\qquad
 C_z=\sum_{w\leq z}p_{w,z}T_w.
\]
The bar involution sends $v$ to $v^{-1}$ and $T_s$ to $T_s^{-1}$.
Under boundedness the source's $a$-function can be described as
\[
 \mathbf a(z)=\max_{x,y}\deg_v h_{x,y,z}.
\]
The maximum is finite; zero polynomials do not contribute.
The leading coefficients are indexed by
\[
 h_{x,y,z}=\gamma_{x,y,z^{-1}}v^{\mathbf a(z)}
                  +\text{terms of smaller degree}.
\]
The coefficient at that power is allowed to be zero for a particular
pair $x,y$. Also write
\[
 p_{1,z}=n_zv^{-\Delta(z)}+\text{terms of smaller degree},
 \qquad\mathcal D=\{z:\mathbf a(z)=\Delta(z)\}.
\]
The left preorder is generated by occurrence of $C_z$ in $C_sC_w$;
right and two-sided preorders are defined analogously. We keep the
source's direction: a term occurring after multiplication is below
the original basis element in this preorder. The cells are the
equivalence classes of the respective preorders.

\paragraph{A complete unequal-weight family.}
Assume every pair of distinct simple generators commutes. Put
$S=\{s_1,\ldots,s_r\}$ and choose arbitrary positive integral weights
$\lambda_i=L(s_i)$. Then $W\cong(\mathbb Z/2)^r$.
Identify an element with its support $I\subseteq\{1,\ldots,r\}$;
write $w_I=\prod_{i\in I}s_i$, and put
$\lambda(I)=\sum_{i\in I}\lambda_i$. Thus inversion fixes every
$w_I$, multiplication in the group is symmetric difference, and Bruhat
order is inclusion of supports.

For one generator of weight $\lambda$, the quadratic relation gives
\[
 T_s^2=1+(v^\lambda-v^{-\lambda})T_s,
 \qquad C_s=T_s+v^{-\lambda},
 \qquad C_s^2=(v^\lambda+v^{-\lambda})C_s.
 \tag{464.1}
\]
The middle expression is bar-invariant: substituting
$T_s^{-1}=T_s-(v^\lambda-v^{-\lambda})$ proves it directly.
Commutativity and uniqueness of the triangular bar-invariant basis give
\[
 C_{w_I}=\prod_{i\in I}(T_{s_i}+v^{-\lambda_i})
       =\sum_{J\subseteq I}v^{-\lambda(I\setminus J)}T_{w_J}.
 \tag{464.2}
\]
Every proper $J$ has a strictly negative exponent because all weights
are positive, as required by the normalization.

It is useful to check the assumed boundedness independently of the
conjectures. In the standard basis,
\[
 T_{w_I}T_{w_J}
 =T_{w_{I\mathbin\triangle J}}
   \prod_{i\in I\cap J}\bigl(1+(v^{\lambda_i}-v^{-\lambda_i})T_{s_i}\bigr).
\]
The supports in each resulting term are distinct. Every coefficient
has degree at most $\lambda(I\cap J)\leq\lambda(S)$, so the
precise standard-basis boundedness hypothesis holds with bound
$\lambda(S)$. There is no appeal to P1--P15 in obtaining this bound.

\paragraph{All structure constants and leading data.}
Repeatedly applying (464.1) gives the especially simple formula
\[
 C_{w_I}C_{w_J}
   =d_{I\cap J}(v)C_{w_{I\cup J}},\qquad
 d_U(v)=\prod_{i\in U}(v^{\lambda_i}+v^{-\lambda_i}),\quad d_\varnothing=1.
 \tag{464.3}
\]
Thus $h_{I,J,K}$ is zero unless $K=I\cup J$, in which case it is
$d_{I\cap J}$. Its degree is $\lambda(I\cap J)$ with leading
coefficient one. For fixed $K$, this degree is at most $\lambda(K)$,
and the upper value is attained by $I=J=K$. Therefore
\[
 \mathbf a(w_K)=\lambda(K).
 \tag{464.4}
\]
Equation (464.2) gives $p_{1,w_K}=v^{-\lambda(K)}$, so
\[
 \Delta(w_K)=\lambda(K),\quad n_{w_K}=1,\quad\mathcal D=W.
 \tag{464.5}
\]
Finally, a nonzero leading coefficient at degree $\mathbf a(w_K)$
requires $I\cup J=K$ and $\lambda(I\cap J)=\lambda(K)$.
Strict positivity of all weights forces $I=J=K$. Since inversion is
trivial here, the source's indexing convention yields
\[
 \gamma_{w_I,w_J,w_K}=
 \begin{cases}1,&I=J=K,\\0,&\text{otherwise}.\end{cases}
 \tag{464.6}
\]
This includes the identity element: its $a$-value and $\Delta$ are
zero, and its only nonzero leading triple is the all-identity triple.

\paragraph{Cells and their order, without an appeal to conjectures.}
Multiplying $C_{w_J}$ on the left by $C_{s_i}$ either leaves its support
unchanged, with scalar $v^{\lambda_i}+v^{-\lambda_i}$, or adjoins
$i$ to the support. Hence a sequence of elementary left-preorder
relations can only add indices, and any desired added indices can
be adjoined one at a time. Consequently
\[
 w_I\leq_L w_J\quad\Longleftrightarrow\quad I\supseteq J.
 \tag{464.7}
\]
The algebra is commutative, so this also describes the right and
two-sided preorders. All left, right and two-sided cells are singletons.
In particular, a larger support has a larger $a$-value, consistent
with the direction of P4. Reversing the direction of (464.7) would
give the wrong sign in that property.

\paragraph{An explicit check of P1--P14.}
All indices below refer to the exact list in [S1, Conjectures 14.2].
We give the verification rather than replace the list by a claim that
the cells look plausible.

\begin{itemize}
\item P1 follows from the equality $\mathbf a=\Delta$ in (464.4)--(464.5).
For P2, a nonzero distinguished leading triple forces all three
elements equal by (464.6), and every element is its own inverse.
For P3 the unique distinguished element associated to $y$ is $y$ itself.
\item For P4, reverse inclusion in (464.7) gives the required
nondecrease of $\lambda(I)$. P5 holds because both the relevant
leading coefficient and $n_d$ are one. P6 holds because $W$ is an
elementary abelian group of exponent two.
\item P7 is immediate from the symmetric diagonal expression (464.6).
For P8 a nonzero triple has three identical elements; their inverses
are the same elements, so all required left-cell relations hold.
\item In P9, P10 and P11, the preorder first gives $I\supseteq J$.
Equality of $a$-values then gives
$\sum_{i\in I\setminus J}\lambda_i=0$, whence $I=J$.
Thus the corresponding cells agree, as required.
\item For P12, a parabolic subgroup generated by $U\subseteq S$
is the same commuting construction on the smaller index set $U$.
The calculation of $\mathbf a(w_I)$ for $I\subseteq U$ is still
$\lambda(I)$, whether made inside the subgroup or in $W$.
\item For P13, each left cell is $\{w_I\}$ and contains the unique
distinguished element $w_I$; the required leading coefficient is one.
For P14, inversion fixes the element, so it fixes its two-sided cell.
\end{itemize}

The positivity of the weights is used essentially in P9--P11.
Equality of sums alone does not mean equality of supports: for weights
$1,2,3$, the supports $\{3\}$ and $\{1,2\}$ have equal sums.
The proof first needs comparability by inclusion. This point also
matters in checking P15.

\paragraph{P15 with two genuinely independent variables.}
Let $v'$ be independent of $v$, and put
$h'_{a,b,c}=h_{a,b,c}(v')$. For the required quadruple write
$w=w_W$, $y=w_Y$, $x=w_X$, $x'=w_{X'}$ and assume
$\mathbf a(w)=\mathbf a(y)$, equivalently $\lambda(W)=\lambda(Y)$.
The identity to verify is
\[
 \sum_{J}h'_{W,X',J}\,h_{X,J,Y}
 =\sum_{J}h_{X,W,J}\,h'_{J,X',Y}.
 \tag{464.8}
\]
On the left a nonzero term requires $J=W\cup X'$ and
$Y=X\cup J=X\cup W\cup X'$. The right has a nonzero term
only under exactly the same final union condition, with
$J=X\cup W$. If that condition fails, both sides vanish.

If it holds, then $W\subseteq Y$. The assumed equality of positive
weight sums forces $Y=W$, and consequently $X,X'\subseteq W$.
Both possible intermediate supports are then $W$. Each side of
(464.8) is exactly
\[
 d_X(v)\,d_{X'}(v').
\]
This proves the two-variable identity in the full Laurent ring
$\mathbb Z[v^{\pm1},(v')^{\pm1}]$, not merely after setting
$v'=v$. It completes all fifteen checks for the commuting family.

The equal-$a$ hypothesis cannot be suppressed even here. In rank one
take $w=1$, $y=s$, and $x=x'=s$. The left side of (464.8) is
$v^{\lambda}+v^{-\lambda}$ while the right side is
$(v')^{\lambda}+(v')^{-\lambda}$. Their specializations at $v'=v$
agree, but the two-variable polynomials do not. The quadruple is
excluded by P15 because $\mathbf a(1)=0\ne\lambda=\mathbf a(s)$.
This is an exact diagnostic against confusing P15 with ordinary
associativity at a single parameter.

\paragraph{What the leading ring looks like in this subcase.}
Using the source's asymptotic multiplication on symbols $t_w$,
\[
 t_xt_y=\sum_z\gamma_{x,y,z^{-1}}t_z,
\]
equation (464.6) gives orthogonal idempotents:
$t_{w_I}t_{w_J}=0$ for $I\ne J$ and $t_{w_I}^2=t_{w_I}$.
The resulting ring is $\prod_{I\subseteq S}\mathbb Z$, with identity
$\sum_I t_{w_I}$. This is a consequence of the explicit calculation,
not an assumption of semisimplicity used to prove it.

\paragraph{The first braid interaction breaks the Boolean model.}
Consider type $A_2$, with $sts=tst$ and the common weight $a>0$.
The weights must be equal: length-additivity gives
$2L(s)+L(t)=L(sts)=L(tst)=L(s)+2L(t)$.
Thus treating two arbitrary unequal numbers as weights in this
example would violate the hypotheses before any Hecke calculation.

Put $q=v^a$. Direct multiplication in its standard basis gives
\[
 C_sC_tC_s=C_{sts}+C_s,
\]
where
\[
 C_{sts}=T_{sts}+q^{-1}(T_{st}+T_{ts})
                 +q^{-2}(T_s+T_t)+q^{-3}.
\]
The displayed element is triangular and bar-invariant because it is
$C_sC_tC_s-C_s$, so uniqueness verifies the identification.
The additional term $C_s$ is absent from (464.3). Supports alone no
longer encode products or cells, and incomparable competing terms must
be handled. A formula obtained by replacing every product by a union
therefore fails in the smallest noncommuting finite example.

\paragraph{What a finite computation can certify.}
For the commuting groups the accompanying script computes the standard
products and the KL products independently in integer Laurent
polynomials, extracts degrees and leading coefficients, and checks
the full two-variable P15 sums for the tested positive weight vectors.
The symbolic proofs above cover every rank and every positive weight;
the finite checks validate indexing and signs, especially the excluded
unequal-$a$ quadruple.

For a general infinite Coxeter group, enumerating elements through
length $N$ gives only lower bounds
\[
 a_N(z)=\max_{\ell(x),\ell(y)\leq N}\deg h_{x,y,z}
 \ \leq\ \mathbf a(z).
\]
The existence of a global bound does not tell us at what finite length
the maximum for each $z$ is attained. A finite calculation violating
$\deg h_{x,y,z}\leq\Delta(z)$ would be a valid obstruction to P1,
but agreement in a finite ball would not prove P1 for all elements.
Likewise, constructing leading coefficients using $a_N$ in place of
the true $\mathbf a$ can create spurious distinguished elements and
incorrect cells. One needs an independently justified upper bound
matching an exhibited value, or a structural theorem computing the
actual $a$-function.

\paragraph{Outcome and remaining gap.}
All fifteen properties, including parabolic compatibility and the
unspecialized P15 identity, have been proved for finite products of
rank-one systems with arbitrary positive integral weights. This is an
established elementary subfamily of known results, with an explicit
independent verification here. No claim for arbitrary bounded weighted
Coxeter groups follows from it.

The first missing step in the proposed extension is a structural
control of KL multiplication once noncommuting braid relations create
additional terms, strong enough to determine the true leading degrees
and their cell relations. Boundedness only limits those degrees;
it does not supply that structure or the two-variable compatibility.
The full conjunction in the catalogue remains unresolved by this attempt.

\subsection{Sources}
\begin{itemize}
\item[S1]Hecke algebras with unequal parameters, revised edition.
\url{https://arxiv.org/pdf/math/0208154v2}.
\textit{Formulation source. Consulted 24 September 2026: Section 13.2 boundedness and the fixed list P1–P15 on printed page 58.}
\item[E1]Xun Xie, \emph{Conjectures P1--P15 for Coxeter groups with
complete graph}, Advances in Mathematics 379 (2021), 107565.
\url{https://arxiv.org/abs/1903.00078};
\url{https://doi.org/10.1016/j.aim.2021.107565}.
\textit{Primary abstract inspected 27 September 2026; the theorem and
its right-angled appendix include the commuting subcase proved here.}
\item[E2]Nathan Chapelier-Laget, J\'er\'emie Guilhot, Eloise Little and
James Parkinson, \emph{The asymptotic Plancherel formula and Lusztig's
asymptotic algebra for affine type A}, Selecta Mathematica 32 (2026),
article 90, published 31 August 2026.
\url{https://link.springer.com/article/10.1007/s00029-026-01176-4}.
\textit{Primary introduction inspected 27 September 2026. Its statement
that general unequal-parameter P1--P15 remain conjectural is
distinguished from the equal-parameter affine setting of that paper.}
\end{itemize}
\end{document}
